\documentclass[../Main.tex]{subfiles}

\begin{document}

\section{Evaluation} \label{sec:Evaluation}

Das Ziel des Playtests ist, die acht Designvorschläge (V1--V8), die wir aus unseren Spieleanalysen und der Leveldesign-Recherche hergeleitet und in Kapitel~\ref{sec:Umsetzung} umgesetzt haben, zu überprüfen, also zu prüfen, ob diese die beabsichtigten Wirkungen auf Spielfluss, Schwierigkeit und Spielgefühl tatsächlich erzeugen.

Wir sind beim Playtest wie in Abschnitt~\ref{sec:VorgehenPlaytest} beschrieben vorgegangen: Zuerst wurde die Testperson mit dem Pre-Questionnaire eingeschätzt, danach folgte der Playtest, der aufgenommen wurde, und anschliessend das Post-Questionnaire sowie das semistrukturierte Interview für tiefergehende Fragen.

Das folgende Kapitel beschreibt zunächst den Ablauf der Playtests und geht dann auf die für die Erkenntnisfrage relevanten Ergebnisse ein. Darauf folgt eine Diskussion, in der wir die Designvorschläge einzeln durchgehen, und schliesslich unsere Empfehlungen für weitere Playtests.

\begin{itemize}[leftmargin=1.5cm]
    \item \textcolor{red}{Abschnitt~\ref{sec:Verbesserungen} heisst \enquote{Verbesserungsvorschläge}. Meint ihr Empfehlungen für weitere Playtests oder Verbesserungen am Spiel? Begriff angleichen.}
\end{itemize}

\subsection{Durchführung und Testpersonen} \label{sec:Durchfuehrung}

Die Playtests wurden immer von Rodin Arslan durchgeführt; Zhile Zhu war hin und wieder dabei. Gespielt wurde immer mit Maus und Tastatur. Drei Playtests fanden über Discord mit Bildschirmübertragung statt, zwei vor Ort. Alle Testpersonen spielten dieselbe exportierte Version des Spiels (für Mac oder Windows) und benutzten dabei ihre eigene, gewohnte Hardware.

Während des Spielens wurde kein Smalltalk geführt; in entspannten Phasen haben wir mit den Testpersonen allerdings über das Spiel gesprochen. Hatte eine Testperson eine Frage zum Kampfsystem oder etwas vergessen, wurde ihre Frage beantwortet. Ganz zu Beginn wurde den Testpersonen zudem gezeigt, welche Aktionen es gibt und wie man sie ausführt, da der Prototyp kein Tutorial enthält.

\begin{itemize}[leftmargin=1.5cm]
    \item \textcolor{red}{Ergänzen: Zeitraum der Playtests, ungefähre Dauer pro Testperson, wie aufgenommen wurde (auch bei den Discord-Tests).}
    \item \textcolor{red}{Abweichungen vom Plan kurz benennen: Einführung durch uns statt Tutorial, Gespräche in entspannten Phasen, Remote-Tests über Discord (Beobachtung nur über den Bildschirm). Diese Punkte auch in den Limitationen (Kapitel~\ref{sec:Schlussfolgerungen}) aufgreifen.}
    \item \textcolor{red}{Anonymisierte Tabelle der Testpersonen TP1--TP5: Spielzeit pro Woche (P1), Lieblingsgenres (P2), Selbsteinschätzung (P5), Spielstil (P6); dazu 2--3 Sätze zur Zusammensetzung.}
    \item \textcolor{red}{Die alte Ablaufbeschreibung (Pre-Questionnaire aushändigen usw.) habe ich gelöscht, sie steht bereits in Abschnitt~\ref{sec:MethodePlaytest}, inklusive der Abweichung vom Think-aloud.}
\end{itemize}

\subsection{Ergebnisse}

Die einzelnen Antworten der Testpersonen sind im Anhang aufgeführt (Abschnitt~\ref{sec:PlaytestProtokolle}). Im Folgenden werden die Ergebnisse pro Designvorschlag beschrieben.

\subsubsection{Ergebnisse pro Designvorschlag} \label{sec:ErgVorschlag}

\paragraph{V1: Überlegenheit als Grundzustand}
Der Aussage, sich den Gegnern grundsätzlich überlegen gefühlt zu haben (Q1), stimmten vor allem jene Testpersonen zu, die auch sehr aggressiv gespielt haben. TP3 stimmte beispielsweise völlig zu; diese Testperson ignorierte auch ihre eigenen Lebenspunkte und konzentrierte sich nur auf die Gegner. TP5 spielte dagegen weniger aggressiv, eher taktisch und methodisch, und bewertete die Aussage neutral. TP4 meinte: \enquote{Ich geniesse meine Überlegenheit}; besonders befriedigend sei es, alle Gegner auf einmal zu besiegen. TP5 hätte sich den Spieler im Vergleich zu den Gegnern hingegen noch stärker vorgestellt, etwa wie in \game{Hyrule Warriors}. Insgesamt gaben drei von vier Testpersonen an, sich überlegen gefühlt zu haben.

Die Überlegenheit wurde aber auch immer wieder durchbrochen. Insbesondere wenn die Testpersonen von einer Gegnergruppe umzingelt waren, konnte es passieren, dass sie wie in einem Stunlock gefangen waren und erst nach erlittenem Schaden wieder herauskamen. In der Luft waren sie zudem durch die Archers verwundbar: Wurden sie getroffen, fielen sie zu Boden, und es kam vor, dass sie beim Aufstehen weiteren Schaden erlitten. Auch die Trägheit der Combos machte die Testpersonen anfällig. Während einer Combo spielte dies eine geringere Rolle, weil die Gegner den Spieler dann kaum angreifen konnten; am Ende einer Combo konnte es jedoch passieren, dass man Schaden erlitt. Kam während einer Combo ein Pfeil angeflogen, konnte man fast nicht ausweichen, und viele Testpersonen haben den Schaden einfach in Kauf genommen.

\begin{itemize}[leftmargin=1.5cm]
    \item \textcolor{red}{Verweise auf den Anhang ergänzen, z.\,B. TP3 (B1), TP4 (B11), TP5 (Q14), Stunlock (TP5 Q13, TP4 B2), Treffer in der Luft (TP2 I6, TP4 B11, TP5 I8), Combos (TP3 I6). Gilt für alle Vorschläge.}
    \item \textcolor{red}{\enquote{Während einer Combo \ldots{} weil die Gegner den Spieler dann kaum angreifen konnten}: Euer Satz war hier unklar (\enquote{weil nicht einen angreifen konnte}). Prüfen, ob es so gemeint ist.}
    \item \textcolor{red}{\enquote{viele Testpersonen} $\rightarrow$ konkret: drei von vier tankten Treffer (TP2, TP3, TP4; B4).}
    \item \textcolor{red}{Fehlt: Der Heavy Soldier, der die Überlegenheit gezielt brechen soll, wurde als eher harmlos erlebt (TP4 I2: verfehlt die Hälfte der Zeit; TP3 I5 und B3; TP2 B1 tankt ihn). Wichtig für die Diskussion.}
\end{itemize}

\paragraph{V2: Begrenzung gleichzeitiger Angreifer}
Alle vier Testpersonen stimmten der Aussage, auch bei vielen Gegnern die Kontrolle behalten zu haben (Q2), eher oder völlig zu. Sie hatten dabei alle eine eigene Strategie, um mit grösseren Gegnermassen umzugehen. Einerseits liessen sich ähnliche Muster beobachten: Einige griffen einfach den ersten Gegner an, der bei ihnen war, andere wollten zuerst alle Archers loswerden. Andererseits gab es unterschiedliche Ausprägungen: TP2 rannte weg, um sich heilen zu können, während TP5 auf die Boxen sprang, die er gleich zu Beginn eines Encounters von Gegnern befreit hatte, um einen sicheren Rückzugsort zu haben. TP3 kämpfte immer gegen viele Gegner gleichzeitig; dabei kam es vor, dass die Testperson einen kurzen Moment der Panik hatte, wenn sie zu wenig Leben hatte. TP5 versuchte eher, Gruppen zu meiden und nur einzelne Gegner anzugreifen.

TP2 meinte, es sei schwerer, wenn die Gegner in einem Haufen seien, als wenn sie einzeln kämen, aber nicht unangenehm schwer. TP5 meinte, er könne nicht einschätzen, wie viel Schaden eine Gruppe von Gegnern mache, wenn er in sie hineinlaufe, und ob er dann im Haufen auch noch in einen Stunlock gerate \textcolor{red}{[Satz endete mit \enquote{und deswegen} -- vervollständigen]}.

\begin{itemize}[leftmargin=1.5cm]
    \item \textcolor{red}{Korrigiert: In eurem Text stand zweimal \enquote{TP2}. Das Säubern der Boxen als Rückzugsort war laut Anhang TP5 (B3), bitte prüfen.}
    \item \textcolor{red}{Fehlt: TP4 häufte Soldiers beim Archer an, um alle auf einmal zu treffen (B2).}
    \item \textcolor{red}{Der Satz zu den Mustern (erster Gegner / zuerst Archers) gehört eher zu V3; hier ggf. kürzen.}
\end{itemize}

\paragraph{V3: Ein taktisches Dilemma durch zwei Gegnerrollen}
Im Interview gaben alle Testpersonen an, zuerst die Archers anzugreifen, wenn sich die Möglichkeit bietet. TP2 meinte, sie sei vor allem auf die Chain und Heavy Archers losgegangen und habe sich die Heavy Soldiers für den Schluss aufbewahrt, da die Archers mit höherer Wahrscheinlichkeit trafen und die anderen Gegner besser kontrollierbar waren. TP3 meinte, sie müsse nicht unbedingt auf die Soldiers losgehen, da diese einen ohnehin suchen und auf den Spieler zurennen. TP4 hat das taktische Dilemma zwar gesehen, meinte aber, es habe für ihn keine allzu grosse Rolle gespielt; er habe einfach, wie er wollte, alles kurz und klein gehauen. Erst wenn er merkte, dass es so nicht ging, und dazu gezwungen war, habe er sich Gedanken zum Vorgehen gemacht.

Beobachtet wurde vor allem bei TP2 und TP3 das erwartete Verhalten: zuerst die Archers, dann die Soldiers, wobei die Soldiers den Weg zu den Archers blockierten. TP4 schlug dagegen zuerst auf die Soldiers ein und ging erst danach auf die Archers los. TP5 probierte Verschiedenes aus: Manchmal griff er zuerst die Archers an, manchmal die Nahkämpfer. Ausschlaggebend waren für ihn eher die Gegner auf den Boxen, weil er sich so einen sicheren Platz zum Heilen verschaffen konnte.

\begin{itemize}[leftmargin=1.5cm]
    \item \textcolor{red}{Angepasst: Euer erster Satz lautete, das Dilemma sei von allen \emph{wahrgenommen und beobachtet} worden, dass alle zuerst auf die Archers gehen. Im Interview sagen das alle (I5), beobachtet wurde es aber nur bei TP2 und TP3 klar. Q3 als Tendenz ergänzen (TP4 neutral, die anderen stimmen zu).}
    \item \textcolor{red}{TP2 laut Anhang: zuerst Archers mit drei Pfeilen oder zielsuchenden Pfeilen; die \enquote{mit den Hämmern} zum Schluss, weil langsam und konterbar (I5). \enquote{Ignorierte} habe ich daher zu \enquote{für den Schluss aufbewahrt} geändert.}
\end{itemize}

\paragraph{V4: Klar telegrafierte Angriffe}
Bei der Aussage, die Angriffe seien klar angekündigt worden (Q4), gehen die Meinungen am weitesten auseinander. Interessanterweise stimmte TP3, die sehr aggressiv und ohne viel nachzudenken spielte, völlig zu. TP2, die ebenfalls sehr aggressiv spielte und meistens erst reagierte, nachdem ein Treffer gelandet war, bewertete die Aussage neutral. Auch TP5, der eher taktisch spielte, bewertete sie neutral, da er die Ankündigungsanimationen der Archers verwirrend fand. TP4, der Projektilen auswich, wenn er konnte, sonst aber alle Treffer in Kauf nahm, stimmte überhaupt nicht zu. Dazu meinte er, die Archers seien besonders, da sie ausserhalb des Blickfelds angreifen könnten und es keinen Sound gebe, der den Angriff ankündigt. TP5 meinte zudem, die Tode seien wegen der Ankündigungen nicht unfair gewesen: Man habe kommen sehen, dass man stirbt, und es habe sich eher wie ein selbstverschuldetes Problem angefühlt.

\begin{itemize}[leftmargin=1.5cm]
    \item \textcolor{red}{Verschoben: Eure Erklärungen (\enquote{das kann damit zusammenhängen \ldots}, \enquote{das könnte darauf hindeuten \ldots}) sind Interpretationen und stehen jetzt in Abschnitt~\ref{sec:Interpretation} unter V4.}
    \item \textcolor{red}{\enquote{ausserhalb des Blickfelds}: Im Anhang steht bei TP4 nur \enquote{ohne wirkliche Prediction oder Sound-Effekte} (I5). Falls das aus der Aufnahme stammt, im Anhang ergänzen.}
    \item \textcolor{red}{Ergänzen: TP3 sah die Pfeile kommen, konnte wegen der Combos aber nicht ausweichen (I6, Q13).}
\end{itemize}

\paragraph{V5: Lesbare Platzierung, dynamisches Verhalten}
hier stimmen tp2 und 3 eher yu und tp 4 un 5 eher nicht. tp2 klettert ywischen wellen auf boxen um den raum zu überblicken und tp3 greif einfach das an was vor ihm steht also reicht ihm die information aus die er hat. in raum 10 der eine enge gasse ist mit winkeln, findet tp3 als einziger den versteckten archer ne zeit lang nicht. tp5 greift auch eher das an was vor ihm liegt und versteht das dynamische verhalten der gegner. er meint einfach dass man, wenn er in den raum geht nicht weiss wo die gegner spawnen bevor man nicht in den raum gegangen ist. tp5 kreist auch hin und wieder eine weile im raum bevor er angreift um die gegner zu beobachten und da er so nicht getroffen werden kann. er sammelt auch so die gegner an. tp4 findet es nicht angenehm wen er von hinten also ausserhalb vom sichtfelt getroffen werden kann. tp4 schaut sich auch eher lang die welle an bevor er angreift. tp3 merkt dass genger ausserhalb des sichtfelds spawnen aber meint dass man sie früh genug bemerke umd reagieren zu können. und er merkt auch dass die archers eher erhöht spawnen und dass es definitiv eine händisch platzierte konstantte geben muss. Also man sieht definitiv muster. tp2 meint ihr gefiel die der gegner im raum gefiele ihr am besten. tp5 meinte nach dem sterben, dass es richtig cool anfühlt wenn man weiss wo der gegner spawnen wird und ihn dierekt zu begin ausnehmen kann.
\begin{itemize}[leftmargin=1.5cm]
    \item \textcolor{red}{Noch zu schreiben. Q5: TP2 und TP3 stimmen zu, TP4 und TP5 eher nicht.}
    \item \textcolor{red}{Beobachtung (B6): TP2 klettert zwischen Wellen auf Boxen, um den Raum zu überblicken; TP3 und TP5 greifen an, was vor ihnen ist; TP3 findet einen versteckten Heavy Archer nicht; TP4 schaut sich Wellen zuerst an, stirbt aber in Raum~9, weil er einen Archer übersieht; TP5 wird fast in einen Türrahmen gedrängt.}
    \item \textcolor{red}{Aussagen: TP3 -- Gegner spawnen ausserhalb der Sicht, werden aber früh genug bemerkt; Archers erhöht und in Gruppen; erste Welle gesetzt (I8). TP2 -- die Verteilung der Gegner gefiel am besten (Q12). TP5 -- am besten, wenn man weiss, wo die Gegner spawnen (Q12). TP4 -- von hinten angegriffen werden ist nicht cool; Gegner vorhersagen macht Spass (B11).}
\end{itemize}

\paragraph{V6: Der Raum als Druckmittel}
Hier gehen die Meinungen etwas auseinander. TP2 spielte in grösseren und kleineren Räumen unterschiedlich: In grossen Räumen spielte sie aggressiver, in kleineren weniger aggressiv und eher taktisch, da es dort wenig Platz zum Heilen gab und das Leben dadurch wertvoller war. Bei TP3 war es unterschiedlich: In manchen kleinen Räumen spielte sie genau gleich aggressiv, in anderen hektischer. TP5 bemerkte dagegen einen deutlichen Unterschied in kleinen Räumen. Wenn diese verwinkelt waren, empfand er sie als gruselig und schaute deshalb vorsichtig um die Ecken.

TP2 meinte, es gebe eine gute Mischung der Raumgrössen. Die kleinen Räume seien tendenziell schwerer, weil es keine Zufluchtsmöglichkeit gebe; ihr Lieblingsraum sei aber Raum~6 mit der Brücke gewesen. TP3 meinte, kleine Räume seien einfacher zu überblicken gewesen, man habe direkt gewusst, was gespawnt sei und was auf einen zukomme. In grösseren Räumen mit Boxen und Brücke habe es dagegen mehr Möglichkeiten gegeben, sich im Raum zu bewegen. Für TP4 spielte die Raumgrösse keine Rolle, weil er einfach alles kopfvoran versucht hat. TP5 meinte, kleine Räume seien übersichtlicher, man könne aber in Ecken gedrängt werden; die grossen seien schwerer gewesen, weil es dort mehr Gegner gab. Wie bereits bei V2 beschrieben, hatte TP5 etwas Mühe mit grösseren Gegnermengen.

\begin{itemize}[leftmargin=1.5cm]
    \item \textcolor{red}{Korrigiert: \enquote{für TP2 spielte es keine Rolle} war laut Anhang TP4 (I4), bitte prüfen.}
    \item \textcolor{red}{Verschoben: Eure Erklärung zu TP3 (\enquote{das liegt wahrscheinlich daran \ldots}) steht jetzt in Abschnitt~\ref{sec:Interpretation} unter V6.}
    \item \textcolor{red}{Q6 als Tendenz ergänzen: Die Antworten gehen weit auseinander (TP2 und TP5 stimmen völlig zu, TP3 neutral, TP4 eher nicht).}
    \item \textcolor{red}{Raum~6 mit der Brücke: Im Anhang nennt TP2 \enquote{Raum 7.5}. Einheitlich machen. \enquote{gruselig} (TP5) steht nicht im Anhang -- falls aus der Aufnahme, dort ergänzen.}
\end{itemize}

\paragraph{V7: Wellenstruktur mit Erholungsphasen}
Alle vier Testpersonen organisierten ihre Erholungszeit selbst. TP2 machte Pause, bevor sie in den nächsten Raum ging, und nutzte die Boxen. TP3 ist gestorben, weil sie sich anfangs nicht geheilt hat; nachdem sie das bemerkt hatte, legte sie am Ende eines Raumes eine Ruhephase zum Heilen ein. Auch TP4 heilte sich vor dem Raumwechsel. TP5 ging dabei sogar etwas taktischer vor und liess oft einen Gegner am Ende einer Welle am Leben, um sich vor der nächsten Welle zu heilen, da man nicht genau wusste, was danach kommt.

TP2 meinte, es sei schwer, während eines Encounters Heilpausen zu finden. TP3 meinte, es gebe eine gute Balance; sobald alle Fernkämpfer besiegt seien, werde es entspannt. TP4 empfand den Verlauf der Kampfabschnitte ebenfalls als angenehm: Er sei immer auf das Spiel konzentriert gewesen, und nichts habe seine Vertiefung gestört. TP5 meinte, die Spannung hänge vom eigenen Leben ab, da man mit wenig Leben angespannter sei.

\begin{itemize}[leftmargin=1.5cm]
    \item \textcolor{red}{Verschoben: \enquote{Das liegt daran, dass TP2 die Wellenstruktur nicht erkannt hat} ist eine Interpretation und steht jetzt in Abschnitt~\ref{sec:Interpretation} unter V7.}
    \item \textcolor{red}{TP5 \enquote{liess oft einen Gegner am Leben}: Im Anhang steht \enquote{heilt sich, wenn am Ende einer Welle nur noch wenige Gegner übrig sind} (B8). Prüfen und angleichen.}
    \item \textcolor{red}{Q7 als Tendenz ergänzen (TP3 und TP5 stimmen völlig zu, TP2 und TP4 neutral). TP2 sprach zudem von \enquote{überlappenden Wellen} (I7) -- wichtig für die Diskussion.}
\end{itemize}

\paragraph{V8: Gestaffelte Einführung neuer Gegner}
Hier stimmten alle Testpersonen völlig zu (Q8). Alle erkannten neue Gegnertypen bei ihrer Einführung an der Farbe, am Aussehen oder an ihren Angriffen. Wie sie mit ihnen umgingen, war allerdings unterschiedlich: TP2 hielt meist kurz inne und analysierte, ob sich noch etwas im Raum befand. TP4 beobachtete den Chain Archer zuerst, ohne anzugreifen, um zu sehen, was er macht. TP3 und TP5 griffen neue Gegner sofort an. TP3 dachte bei der Einführung des Archers, es handle sich einfach um eine stärkere Version der Soldiers, und wusste anhand des Aussehens anfangs nicht, was die neuen Gegner tun.

\begin{itemize}[leftmargin=1.5cm]
    \item \textcolor{red}{Allgemein für 6.2.1: Pronomen einheitlich halten (ich habe für TP2 und TP3 \enquote{sie} wie im Anhang verwendet, für TP4 und TP5 \enquote{er}). Für die Anonymisierung evtl.\ durchgehend \enquote{die Testperson}.}
    \item \textcolor{red}{Sobald TP1 da ist: \enquote{x von 4} zu \enquote{x von 5} anpassen.}
\end{itemize}

\subsubsection{Gesamteindruck und allgemeine Beobachtungen} \label{sec:ErgAllgemein}

Der gesamteindruck aller spieler zu unseren 3 wertungen fällt ziemlich ähnlich aus. für das spielgefühl meinten alle es sei eher befiredigend zu spielen. den spielfluss fanden auch alle zusammenhängend dabei meinten sie, dass die Räume zwar alle für sich seperat sind und sie nicht viel verbindet aber der instesammte spielfluss gepasst hätte. Die schwierigkeit haben auch alle anggemessen gefunden, alle sind eins bis zwei mal gestorben aber nicht wiederholt an gleicher stelle. es gab auch viele momente wo sie knapp überlebt haben.

Weiteres was gesagt wurde war folgendes:
was sich am besten angefühlt hat waren gegnerverteilung und Spawns, wenn man gegner spawnkillen kann. auch die taktischen entscheidungen und heilen und weglaufen wären interessant. Horden zu vernichten, was der sinn des spieles ist, hätet auch hier spass gemacht und sich gut angefühlt. die tiefe, die wir mit unterschiedlichen combos also light und heavy und deren vershciedene schadenswerte wären auch angenehm.
unfair hätte sich angefühlt, dass man combos nicht abbrechen konnte und auf gegner deshalb nicht so einfach reagieren konnte. wenn man in einer horde von gegnern eingeklämmt fühlte sich das auch unbefriedigend an. jemand meinte auch dass sich das spring verhalten nicht wie erwartet anfühlen würde und das kreutzt sich mit einer anderen ausssage, dass der übergang von movement und combat eetwas unstimmig wäre.
Im vergleich mit den erwartungen waren die playtester positiv überrascht vom umfang und der gegnervielfalt. es wurde ein boss erwartet der nicht ekommen ist, und auch soweas wie block oder doge wurde sich hin und wieder gewünscht, aber es gab auch stimmmen die ihre meinung am ende geändert habe und meinten, dass diese moves nicht reingepasst hätten. man hat auch erwartet das man als spieler noch stärker wäre.

Weitere bemerkungen sind, dass die orientierung nicht immer glatt lief, oft wurde nicht bemerkt dass der encounter fertig ist und offene türen wurden so übersehen, da es eine audiovisuellen indikatoren gab, dass es der fall ist. Bei der steuerung war den spielterstern nicht immer ganz die tiefe der inputs natürlich. also man könnte dash gedrückt halten, spieler die das schon kannten waren entzückt, während die die das nicht kannten eher gespamt haben. springen war ähnlich mit gedrück halten. die reichweite der eigenen waffe einzuschätzen war auch etwas schwer. die dauer der heilung war auch nicht ganz schlüssig. Tp2 meinte es wäre stat mündlich auch ein tutorial als verbesserungsvorschlag mölglich. tp4 und 5 haben versucht mit objekten zu interagieren und suchen beide glitches
die stimmung in allen playtests war positiv tp3 und 4 haben viel gelacht. tp3 und 5 haben das spiel nach dem beenden einfach weiter gespielt. tp2 war hin und wieder auch genervt, dass sie beim heilen unterbrochen wurde. jeder spieler hatte oft eienn gegnertypen der sie besonders genervt hat. tp5 meinte er wäre nie unter stress gewesen.

\begin{itemize}[leftmargin=1.5cm]
    \item \textcolor{red}{Q9--Q11 als Tendenz: Spielgefühl (Q9) bei allen eher positiv; Spielfluss (Q10) überwiegend positiv, TP3 unterscheidet aber \enquote{von Raum zu Raum} (tief) und \enquote{im Kampf} (hoch); Schwierigkeit (Q11) überwiegend angemessen.}
    \item \textcolor{red}{Offene Antworten Q12--Q14 und allgemeine Beobachtungen B10--B14: offene Türen nicht bemerkt (TP2, TP4), fehlendes Tutorial (TP2), Combos nicht abbrechbar und Wunsch nach Block (TP4, TP5), Suche nach Bugs (TP3, TP5), Boss erwartet (TP2, TP3).}
\end{itemize}

\subsection{Diskussion}

\subsubsection{Interpretation pro Designvorschlag} \label{sec:Interpretation}

\begin{itemize}[leftmargin=1.5cm]
    \item \textcolor{red}{Pro V1--V8: bestätigt / teilweise bestätigt / nicht bestätigt, mit Begründung aus den Daten und Bezug zur Recherche. Abweichungen zwischen Vorschlag und Umsetzung (z.\,B. V2, V7) hier diskutieren. Selbst schreiben.}
    \item \textcolor{red}{Unten stehen eure Erklärungen, die ich aus 6.2.1 hierher verschoben habe. Sie sind ein guter Anfang für V4, V6 und V7.}
\end{itemize}

\paragraph{V4: Klar telegrafierte Angriffe}
Dass TP3 völlig zustimmte, kann damit zusammenhängen, dass ihr die meisten Treffer egal waren und ihr die Ankündigungen deshalb genügten. Dass TP2 erst nach einem Treffer reagierte und nur neutral bewertete, könnte darauf hindeuten, dass die Ankündigungszeit etwas zu kurz war. Dass TP4 überhaupt nicht zustimmte, kann damit zusammenhängen, dass er die Information nicht brauchte und deshalb nicht darauf achtete, oder auch daran liegen, dass die Combos zu träge sind und er deshalb auf die Ankündigungen nicht reagieren konnte.

\paragraph{V6: Der Raum als Druckmittel}
Dass TP3 in kleinen Räumen teils gleich aggressiv, teils hektischer spielte, liegt wahrscheinlich daran, dass sie wegen ihres aggressiven Spielstils ohnehin wenig Platz brauchte und den zusätzlichen Platz der Prospect Spaces nur in Momenten mit wenig Leben oder zu grosser Gefahr vermisste.

\paragraph{V7: Wellenstruktur mit Erholungsphasen}
Dass TP2 es schwer fand, Heilpausen zu finden, liegt daran, dass sie die Wellenstruktur des Spiels nicht erkannt hat und deshalb nicht davon profitierte. \textcolor{red}{[Belegt? Woran erkennt man, dass sie die Wellenstruktur nicht erkannt hat? Sonst vorsichtiger formulieren: \enquote{könnte daran liegen}.]}

\subsubsection{Gesamteindruck: Spielgefühl, Spielfluss und Schwierigkeit} \label{sec:Gesamteindruck}

\begin{itemize}[leftmargin=1.5cm]
    \item \textcolor{red}{Auswertung von Q9--Q11 und B14 zusammen mit den Ergebnissen pro Vorschlag. Das ist die Brücke zur Erkenntnisfrage in Kapitel~\ref{sec:Schlussfolgerungen}.}
\end{itemize}

\subsection{Verbesserungsvorschläge} \label{sec:Verbesserungen}

\begin{itemize}[leftmargin=1.5cm]
    \item \textcolor{red}{Konkrete Verbesserungen, jede mit Beleg aus den Daten (z.\,B. akustische Angriffsankündigung, Tutorial, Signal beim Öffnen der Tür, abbrechbare Combos, Block). Laut Vertrag müssen sie nicht mehr umgesetzt werden.}
\end{itemize}

\end{document}
